\section{Conclusion}
\label{sec:conclusion}

This paper introduced CodeGraph, a hybrid retrieval system designed to bridge the Semantic Gap between natural-language task descriptions and repository-scale source code. By extracting lexical seeds and propagating their relevance over a static dependency graph using Personalized PageRank, the system achieves 74.0\% Recall@10 on the SWE-bench Lite benchmark.

The central finding of this work is that structural retrieval performance is primarily bottlenecked by Reachability rather than Ranking. When lexical signals identify even a noisy starting point, the graph topology can reliably arbitrate relevance and push correct files to the top. The Seed-First engineering strategy---incorporating uniform PPR restart, compound tokenization, and IDF edge weighting---yielded a 32 percentage point improvement over a strong lexical baseline.

Future work must address the two persistent failure modes: Semantic Gaps and Connectivity Gaps. Replacing the BM25 lexical stage with a lightweight dense embedding model could map abstract task language to code entities more effectively, solving the semantic gap while still leveraging the graph for structural propagation. For connectivity gaps, integrating runtime call graphs or dynamic plugin registries could expose the implicit dependencies that static analysis currently misses. 

CodeGraph demonstrates that AI agents do not need to rely solely on text search; treating code as a graph provides a powerful, interpretable, and deterministic mechanism for navigating large software repositories.
